\begin{table}[t]
\setlength{\tabcolsep}{3pt}
\centering \footnotesize
\begin{tabular}{r|lll l}
\textbf{Dataset} & \textbf{Audience} & \textbf{Domain} & \textbf{\# Docs} & \textbf{\# Tokens} \\ \hline
% \multicolumn{4}{r}{\textit{Readability Known Datasets}} \\ \hdashline[0.5pt/3pt]
arXiv & Experts & CS &6440 & 163 \\
PubMed & Experts & Medicine &6658 & 205 \\
SciTLDR & Experts & CS &618 & 19 \\
SJK & Kids & Varied & 284  & 142 \\ \hdashline[0.5pt/3pt]
% \multicolumn{4}{r}{\textit{PLS Datasets}} \\ \hdashline[0.5pt/3pt]
CDSR & General & Healthcare &284 & 221 \\
PLOS & General & Biomed &1376 & 195 \\
eLife & General & Biomed &241 & 383 \\
Eureka & Journalists & Varied &1010 & 662 \\
CELLS & General & Biomed &6311 & 162 \\
SciNews & General &  Varied & 4188 & 615 
\end{tabular}
\caption[Comparison of the datasets analyzed in this chapter.]{Comparison of the datasets analyzed in this chapter. The first 4 are datasets with a specific target audience. The following 6 datasets are commonly used in PLS literature. We report the number of documents (\# Docs) in the test set as well as the average number of tokens (\# Tokens).
% \daniel{FYI: change D and T to Docs/Tokens since you haves room.}
}\label{tab:dataset-comp}
\end{table}